% Chunk: Chapter 7 backmatter---Problems and References
% Source: Weinberg QFT Vol. I, physical PDF pp. 360--361 (printed pp. 337--338)
% Figures: none.
% Uncertainties: none.

\chapterbackmatter{Problems}

\begin{enumerate}
\item Consider the theory of a set of real scalar fields $\Phi^n$, with
Lagrangian density
\[
  \mathcal{L}=-\frac12\sum_{mn}\partial_\mu\Phi^n\partial^\mu\Phi^m
  f_{mn}(\Phi),
\]
where $f_{mn}(\Phi)$ is an arbitrary real matrix function of the field. (This
is called the non-linear $\sigma$-model.) Carry out the canonical quantization
of this theory. Derive the interaction
$H_I[\phi(t),\dot\phi(t)]$ in the interaction picture.

\item Consider a theory of real scalar fields $\Phi^n$ and Dirac fields
$\Psi^i$, with Lagrangian density $\mathcal{L}=\mathcal{L}_0+\mathcal{L}_1$,
where $\mathcal{L}_0$ is the usual free-field Lagrangian density, and
$\mathcal{L}_1$ is an interaction term involving $\Phi^n$ and $\Psi^i$,
but not their derivatives. Derive an explicit expression for the symmetric
energy--momentum tensor $\Theta^{\mu\nu}$.

\item In the theory described in Problem 2, suppose that the Lagrangian
density is invariant under a global infinitesimal symmetry
$\delta\Phi^n=\ii\epsilon\sum_m t^n_m\Phi^m$,
$\delta\Psi^i=\ii\epsilon\sum_j\tau^i_j\Psi^j$. Derive an explicit expression
for the conserved current associated with this symmetry.

\item Consider the theory of a complex scalar field $\Phi$ and a real vector
field $V^\mu$, with Lagrangian density
\[
  \mathcal{L}=-(D_\mu\Phi)^\dagger D^\mu\Phi
  -\frac14F_{\mu\nu}F^{\mu\nu}
  -\frac12m^2V_\mu V^\mu-\mathcal{H}(\Phi^\dagger\Phi),
\]
where $D_\mu\equiv\partial_\mu-igV_\mu$ and
$F_{\mu\nu}\equiv\partial_\mu V_\nu-\partial_\nu V_\mu$, and $\mathcal{H}$
is an arbitrary function. Carry out the canonical quantization of this
theory. Derive the interaction in the interaction picture.

\item In the theory of Problem 4, derive expressions for the symmetric
energy--momentum tensor $\Theta^{\mu\nu}$ and for the conserved current
associated with the symmetry under $\delta\Phi=\ii\epsilon\Phi$,
$\delta V^\mu=0$.

\item Prove that the Dirac bracket satisfies the Jacobi identity \eqref{eq:7.6.23}.

\item Prove that the Dirac bracket is independent of the choice of constraint
functions $\chi_N$ used to describe a given submanifold of phase space.
\end{enumerate}

\chapterbackmatter{References}

\begin{enumerate}
\item \label{ref:7.1} H. B. G. Casimir, \emph{Proc. K. Ned. Akad. Wet.} \textbf{51}, 635
(1948); M. J. Spaarnay, \emph{Nature} \textbf{180}, 334 (1957).

\item \label{ref:7.2} F. Belinfante, \emph{Physica} \textbf{6}, 887 (1939); also see
L. Rosenfeld, \emph{M\'emoires de l'Acad\'emie Roy. Belgique} \textbf{6}, 30
(1930).

\item \label{ref:7.3} See, e.g., S. Weinberg, \emph{Gravitation and Cosmology} (Wiley, New
York, 1972): Chapter 12.

\item \label{ref:7.4} See, e.g., J. Wess and J. Bagger, \emph{Supersymmetry and Supergravity}
(Princeton University Press, Princeton, 1983), and original references quoted
therein.

\item \label{ref:7.5} P. A. M. Dirac, \emph{Lectures on Quantum Mechanics} (Yeshiva
University, New York, 1964). Also see P. A. M. Dirac, \emph{Can. J. Math.}
\textbf{2}, 129 (1950); \emph{Proc. Roy. Soc. London}, ser. A, \textbf{246},
326 (1958); P. G. Bergmann, \emph{Helv. Phys. Acta Suppl.} IV, 79 (1956).

\item \label{ref:7.6} T. Maskawa and H. Nakajima, \emph{Prog. Theor. Phys.} \textbf{56},
1295 (1976). I am grateful to J. Feinberg for bringing this reference to my
attention.
\end{enumerate}
